%%%
%
% $Autor: Wings $
% $Datum: 2021-05-14 $
% $Pfad: GitLab/MLEdgeComputer/ArdiuonoBootloader
% $Dateiname: 
% $Version: 4620 $
%
% !TeX spellcheck = de_DE/GB
% !TeX program = pdflatex
% !BIB program = biber/bibtex
% !TeX encoding = utf8
%
%%%




\chapter{Bootloader}\index{Bootloader}

A bootloader is an important component of computer systems and microcontrollers. A bootloader, also known as a boot program or bootstrap loader, is special software that is loaded into the main memory of a computer after startup.
Its main task is to load and execute the actual operating system or application software.

Some microcontrollers have little or no internal flash memory. In such cases, the bootloader takes over the copying of code from non-executable memory (e.g. SPI flash) to an executable area such as RAM. The bootloader enables the separation of critical parts of the program (e.g. firmware updates) from frequently changing application code.
A bootloader can verify a cryptographic signature to ensure that the application has not been replaced or tampered with.

\bigskip

The bootloader must fulfill two main tasks:

\begin{enumerate}
  \item Execute at start of the microcontroller board: The bootloader starts automatically when the microcontroller is switched on.
  \item Jump to application code: After startup, the bootloader loads the application software and jumps to its execution.
\end{enumerate}

\bigskip

\section{Reset of an Arduino}\index{Reset of an Arduino}

During the development of a program, it may happen that the Arduino board is no longer responsive. In this case, a reset must be carried out. With the help of the small button\Mynote{No?} on the Arduino board, it boots into the bootloader and can then be reloaded. In this case, the configuration of the Arduino IDE must be checked. In figure \ref{fig:ArduinoReset} the reset button is shown for the Arduino Nano 33 BLE Sense. \Mynote{tikz picture}

\begin{center}
    \includegraphics[width=0.6\textwidth,angle=90]{Nano33BLESense/reset.png}
    \captionof{figure}{The Reset Button of the Arduino}\label{fig:ArduinoReset}
\end{center}

\section{Bootload for OpenMV}

\URL{https://docs.arduino.cc/tutorials/nano-33-ble/getting-started-omv/}


\section{Kopieren eines Arduinos}

\URL{https://www.instructables.com/How-To-Burn-a-Bootloader-to-Clone-Arduino-Nano-30/}

\bigskip

I have recently purchased an Arduino Nano 3.0 Clone which came without a bootloader. I am sure there are many other people who are in the same situation as me, and may have freaked out a little bit at first ! Don't worry, in this instructable I will show you how to easily install a bootloader to your new clone.

ATTENTION - some clones are different. they don't have an ftdi chip, which may cause your computer to not see the arduino once plugged in with the usb. if this is the case, please follow these steps; https://www.instructables.com/id/How-To-Burn-a-Boot... thanks to author Gegec for his help on noticing the problem and solving it as well :)


\subsection{Step 1: Get a Clone}

It is very easy to find a clone of any board on the internet. Or you can at worst make yourself one...

.

They are extremely cheap, I have seen pinless nano boards as low as \$5 online. You just have to do some work on your end, which is a very fair give and take if you ask me.

.

I don't want to advertise any names here, if you just search the board that you want with the word clone next to it, you will find some. Mine is called Funduino Nano, which is an open source company apparently just making amazing products at real fair prices..

.

Look how small an arduino nano is, it's the size of my pinky finger. This will help me tremendously in my new projects, huge improvements on size!

\subsection{Step 2: Test It First}

Plug it to the computer with the USB cable.

.

Make sure it turns on, and yes it did.! So far we are safe.

.

Now try to upload the blink sketch to it using the Arduino software on your computer. (file>examples>basics>blink)

.

You will most likely get the error in the second picture. This has nothing to do with the code being right/wrong. It just simply means that your computer and the board aren't communicating correctly, in our case the lack of the bootloader prevents us from doing so.

.

If your clone did work with the sketch upload, then even better.! you found one that is pre-installed. just start using it...

If you get the error message, keep on reading..>

\subsection{Step 3: Simple Few Steps}

\begin{center}
  \includegraphics[width=0.6\textwidth,angle=90]{Bootloader/ICPSPinout}
  \captionof{figure}{Pinout of the ICPS}\label{fig:ICPSPinout}\index{ICPSPinout}
\end{center}

\begin{center}
    \includegraphics[width=0.6\textwidth,angle=90]{Bootloader/BootloaderErrorMessage}
    \captionof{figure}{Bootloader Error Message}\label{fig:ICPSPinout}\index{ICPSPinout}
\end{center}

\begin{center}
    \includegraphics[width=0.6\textwidth,angle=90]{Bootloader/PinConnection}
    \captionof{figure}{Connection of two Arduinos}\label{fig:ICPSPinout}\index{PinConnection}
\end{center}





Now you need to find a working arduino board, doesn't matter if it is a nano or not, I have used my ``real'' uno r3 and worked perfectly. If you don't already have one, try borrowing it from a friend.



If you have any grab your female - male breadboard jumpers. These will make the process even faster and easier. BUT I didn't have any so, instead of waiting 2 weeks for them to arrive online as well, I had to take some extreme measures!



First we need to locate and get familiar with the ICSP pins. On mine just one of them was marked with a number, top right as 1. Here is what they look like in the schematics. So right under 1 had to be 2 and next to one was 3 and so on \ldots



Now that you are ready lets begin - do the following steps in order;


\subsubsection{STEP 1 - Upload ISP}

\begin{itemize}
  \item Connect the working arduino to your computer.
  \item Open the arduino software on the computer
  \item Go to Tools>Serial Port>COM (and make sure that the correct COM port is chosen - ask me if you don't what this means)
  \item  Go to Tools>Board>Arduino Uno (in my case uno, this is the arduino that is already working not yet the other one)
  \item Go to File>Examples>Arduino ISP (a new sketch screen will open with the ISP code written in it)
  \item Verify then Upload (your board will start blinking for a bit, it means it is loading, wait until done uploading)
  \item Unplug the arduino. Now it has the ISP code in it, all we have left to do is connect them together and run a few commands again \ldots
\end{itemize}

\subsubsection{STEP 2 - Connect 2 arduinos together}

\begin{itemize}
  \item Now connect the two arduinos together, to do this without female-male jumper cables was time consuming, you have to be VERY CAREFUL and take your time bending the wires, other wise you may short your circuits. ask me first ! In the second picture you can see that first I put the ICSP pins through a piece of paper and then carefully wrapped around them the corresponding wires. Make sure they grab on very tightly, use a small flat head screw driver to bend them. Take your time, don't be rough.
  \item Connect them this way;
    \begin{itemize}
      \item pin 1 on clone to D12 on working arduino
      \item pin 2 on clone to 5V on working arduino
      \item pin 3 on clone to D13 on working arduino
      \item pin 4 on clone to D11 on working arduino
      \item pin 5 on clone to D10 on working arduino
      \item pin 6 on clone to GND on working arduino
    \end{itemize}
\end{itemize}

\subsubsection{STEP 3 - Burn Bootloader}

\begin{itemize}
  \item Keeping all the wires intact and unmoved in step 2, connect the working arduino back to the computer with the USB
  \item Go to Tools>Serial Port>COM (and make sure again that the correct COM port is chosen - ask me if you don't what this means)
  \item Go to Tools>Board>Arduino Nano w/Atmega 328 (this time we choose the non working arduino, in my case nano atmega 328)
  \item Go to Tools>Programmer>Arduino as ISP
  \item Go to Tools>Burn Bootloader (this is it, after the LEDs stop flashing your board is ready to use !!!)
\end{itemize}


\bigskip


Do all of the above steps in ORDER !!! no skipping steps

\bigskip

Always make sure that the correct COMM port and correct board are chosen before uploading a sketch to avoid errors \ldots

\bigskip

if you want to learn more or get a second opinion you can check this \HREF{https://sysexit.wordpress.com/2013/02/07/burning-a-bootloader-to-an-arduino-nano-using-another-arduino/}{post} (it's amazing) or just search on google ``how to install a bootloader to arduino clone'' and you will see many articles on this very topic.

\section{to do}

\begin{itemize}
  \item \URL{https://support.arduino.cc/hc/en-us/articles/4841602539164-Burn-the-bootloader-on-UNO-Mega-and-classic-Nano-using-another-Arduino}
  \item \URL{https://forum.arduino.cc/t/how-to-re-flash-firmware-to-an-arduino-33-ble-sense/1053526}
\end{itemize}
